%-------------------------------------------------------------------------------
%	SECTION TITLE
%-------------------------------------------------------------------------------
\cvsection{Publications}


%-------------------------------------------------------------------------------
%	CONTENT
%-------------------------------------------------------------------------------
\begin{itemize}
    \item{\textbf{Strategic fiscal and monetary interactions in the Brazilian economy}

\small{This paper identifies the leadership structure of the game played by monetary and fiscal authorities in the Brazilian economy after the implementation of inflation targeting regime in 1999. A stylized small-scale New Keynesian model augmented with fiscal policy is estimated using Bayesian methods. I assume that monetary and fiscal authorities can act strategically under discretion in a non-cooperative setup and compare three different forms of games: (i) simultaneous move; (ii) fiscal leadership; and (iii) monetary leadership. I find strong empirical support for the hypothesis that the Brazilian fiscal authority acts as a Stackelberg leader. The results obtained can shed some light on the improvement of policy design in the Brazilian economy. \\ \emph{Revista Brasileira de Economia (Brazilian Review of Economics), v. 80, n. 1 (2026): JAN-MAR. DOI: \href{https://doi.org/10.5935/0034-7140.e042026}{10.5935/0034-7140.e042026}.}}}
\end{itemize}
